\documentclass[quick,pro]{debate}
\motion{对知识网红的崇拜让我们离真知更近（正方）}
\begin{document}
\debatetitle
\begin{thesisbox}[一句话立论]崇拜不是终点，是入口——它把抽象的知识变成一个你愿意反复靠近的人，让本来不会开始的人开始了，也让真专家愿意留在公共场域，所以我们离真知更近。\end{thesisbox}
\section{定义与判准（原话）}
\begin{fields}
\field{知识网红}{在网上传播知识、并因此获得大量关注的创作者。
}
\field{崇拜}{明显强于一般关注的推崇与信任，表现为持续追随、优先采信其观点。\textbf{是情感的强度，不是认知的关闭。}
}
\field{真知}{正确而有理据的认识，三层——内容对，知道自己凭什么认为它对，错了能改。
}
\field{判准}{跟“同样在看这些人、但没有崇拜”的自己比，崇拜这一层在接触、正确性、理据、纠错四项上的净额是往前还是往后。\textbf{我方要证明}接触增量大于理据损耗；\textbf{对方要证明}损耗大于增量。
}
\field{两条对称禁令}{知识网红的功劳不能记在崇拜头上；短视频媒介的毛病也不能记在崇拜头上。
}
\field{对方提偏向定义或判准时的 30 秒口径}{①“无条件”“盲目”是对方加的词，词典里只有“尊敬钦佩”，请用双方共认的定义打这场比赛。②对方要求“能独立复述推理链条”，这个标准把中学教材、科普书、通识课全判成负分，它衡量的不是真知，是专业训练。③题目问的是“更近”，不是“到达”。
}
\end{fields}
\section{我方论点}
\begin{longtable}{@{}L{\dimexpr0.049\linewidth-1.50\tabcolsep}L{\dimexpr0.104\linewidth-1.50\tabcolsep}L{\dimexpr0.493\linewidth-1.50\tabcolsep}L{\dimexpr0.353\linewidth-1.50\tabcolsep}@{}}
\toprule \thead{标签} & \thead{一句话} & \thead{核心数据 / 学理 / 案例} & \thead{出处}\\\midrule\endhead
\textbf{门票效应} & 崇拜是把一次性接触变成持续关注的现实机制 & 学理：Katz \& Lazarsfeld 1955 意见领袖。数据：2023 年公民具备科学素质比例 14.14\%，互联网首选渠道 58.3\%（只证渠道转移，\textbf{不可说成崇拜造成}）。案例：汪品先院士入驻 B 站——“有这么多人可以听讲，这是老师最大的享受” & 中国科协第十三次公民科学素质抽样调查（2024-04-16，有效样本 28.9 万）；第一财经\\
\textbf{撤回地址} & 崇拜给错误一个可以被撤回的地址 & 学理：Hardwig 1985——什么都自己验证的人，信念只会更粗糙、更不理性。数据：166 条疫情短视频中 72\% 不含错误信息。案例：AI 合成张文宏卖蛋白棒 1266 件，真人一出面就被戳穿 & Hardwig, \emph{J. Philosophy} 82(7)；Baghdadi et al. 2023, \emph{JAMIA Open}；央视网 2024-12-18\\
\bottomrule\end{longtable}
\section{对方论点 → 一句话反驳}
\begin{longtable}{@{}L{\dimexpr0.208\linewidth-1.33\tabcolsep}L{\dimexpr0.394\linewidth-1.33\tabcolsep}L{\dimexpr0.398\linewidth-1.33\tabcolsep}@{}}
\toprule \thead{对方会说} & \thead{我方一句话回} & \thead{对方追问时}\\\midrule\endhead
用人替证据，信念支点从证据换成人 & 一个不依赖任何人、只信自己验证过的证据的普通人，现实中是谁？ & 问分界线：是资质？四部门通知正在核验资质。是感情？请证明有感情的信任更容易出错\\
崇拜关掉了“自己解释一遍”这个开关 & 对方三份研究，一份讲搜索、一份讲看电视、一份讲全人类，没有一份自变量是崇拜 & 要数据：崇拜者比不崇拜的人更少查证，这句经验命题对方全场没给过任何证据\\
崇拜可被劫持（AI 换脸、带货、越界） & 假信息能冒充一个人，正因为那个人存在，他才能出来说“这不是我” & 不完美不等于不存在；请举一个匿名谣言被“核验资质”治好的例子\\
算法选的是表达能力不是正确性 & 那是平台的问题，按对方自己承认的规矩不能记在崇拜头上 & 算法能制造一次点击，制造不了三年的追更\\
崇拜者只要结论不要过程 & 人类知识的入口，有哪一个不是从结论开始的？ & 我方承认崇拜不产生深度理解，我方主张它产生方向\\
崇拜让人不读原始文献 & 在崇拜发生之前，这些人读原始文献吗？ & 基线是电视 85.5\%、互联网 79.2\%，这就是真实的对照组\\
\bottomrule\end{longtable}
\section{必问与必答}
\begin{points}{必问 （对方不答就点名、重复、不换题）}
\item 今天有哪一条数据的自变量是崇拜？
\item 起点前移，算不算更近？
\item 你方那个“不崇拜但会去查文献”的普通人，是谁？（一辩稿抛出的问题）
\item 崇拜和信任的分界线在哪？
\item 如果大家信的是一句无名的话，谁有资格出来否认它？
\end{points}
\begin{points}{必答 （15 秒正面回答，不许回避）}
\item “崇拜者知道自己为什么相信吗？”→ 多数不知道全部理由，但知道去找谁核对，也知道那个人是谁。
\item “AI 换脸卖出 1266 件，是崇拜的功劳吗？”→ 是造假的罪过；崇拜的功劳是让能出来辟谣的那个人存在。
\item “从入口到深入的转化率是多少？”→ 给不出，对方也给不出损耗率，双方都在机制层。
\item “你方是不是在为盲从辩护？”→ 不是，盲从是对方加的词，我方用的是双方共认的定义。
\end{points}
\section{自由辩战场概要}
\begin{longtable}{@{}L{\dimexpr0.212\linewidth-1.50\tabcolsep}L{\dimexpr0.201\linewidth-1.50\tabcolsep}L{\dimexpr0.240\linewidth-1.50\tabcolsep}L{\dimexpr0.346\linewidth-1.50\tabcolsep}@{}}
\toprule \thead{战场} & \thead{开场问题} & \thead{目标} & \thead{收束语}\\\midrule\endhead
一·归因（前 90 秒） & 今天有哪一条数据的自变量是崇拜？ & 让评委记住对方的证据不是关于崇拜的 & “对方骂的是媒介，寄来的账单写着崇拜。”\\
二·位移（中 90 秒） & 起点前移，算不算更近？ & 把判准从“到达”拉到“位移” & “我方从不说他到了。我方说他动了，对方一步都没还给我们。”\\
三·收束（后 60 秒） & 不开新战场 & 重复一句话立论与必问 & “这笔账要记对人。”\\
\bottomrule\end{longtable}
\section{如果……就……}
\begin{contingency}
\ifthen{对方把崇拜定成“无条件盲从”，或要我方证明崇拜者理解更深}{“这四个字是对方加的，词典里只有‘尊敬钦佩’”；理解深浅不接，只答“我方主张的是位移”。}
\ifthen{对方拿出没预料到的数据}{先问出处与自变量：“这项研究测的自变量是崇拜，还是媒介？”}
\ifthen{门票效应被打穿，或对方回避必问}{收缩到“撤回地址”，用判准第四项解释为什么仍占优；对方回避就点名、重复、不换题。}
\end{contingency}
\end{document}
